%% ==================================================
%% Chapter: Sequential generation of matrix product states
%% Source: arXiv:quant-ph/0608197, section "Generation of MPS"
%% (lines 1504--1616 of Papers/quant-ph_0608197/MPSarchive.tex), and the
%% successive-decomposition construction of arXiv:quant-ph/0501096,
%% eq. induction.
%% ==================================================

\chapter{Sequential Generation of Matrix Product States}\label{ch:seqanc}

A chain of $N$ sites with local dimension $d$ starts in $\ket{0}^{\otimes N}$,
and an ancilla $\C^D$ interacts with the sites one after another. This chapter
records the probabilistic and deterministic schemes of
\cite[Section~5.1]{PerezGarcia2007Matrix} and the characterization of the states
they produce by open-boundary matrix product representations of bond dimension
at most $D$. The last two sections treat the deterministic transition scheme
for $d=2$ and sequential generation without an ancilla.

Sites are labelled $k=0,\dots,N-1$ in the order of the tensor factors of the
ket, and a configuration is $\sigma=(\sigma_0,\dots,\sigma_{N-1})$. The source
writes the generated state as $\sum\bra{\varphi_F}A^{[N]}_{i_N}\cdots
    A^{[1]}_{i_1}\ket{\varphi_I}\ket{i_N\cdots i_1}$; in the present labelling
the matrix $A_k$ at ket position $k$ is the source's $A^{[N-k]}$, the operation
applied at step $N-k$. For square matrices $Q_0,\dots,Q_{n-1}$ indexed by the
physical index write $Q^\sigma=Q_0^{\sigma_0}\cdots Q_{n-1}^{\sigma_{n-1}}$.

\section{Sequential schemes}

\begin{definition}[Sequentially generated state]
    \label{def:seqanc_state}
    \lean{MPSPreparation.seqAmplitude}
    \leanok
    Let $\ket{\varphi_I},\ket{\varphi_F}\in\C^D$ and let
    $A_k^i\in\MN{D}$ for $k=0,\dots,N-1$ and $i=0,\dots,d-1$. The state
    produced by the operations
    $\sum_{i,\alpha,\beta}(A_k^i)_{\alpha\beta}\ketbra{\alpha,i}{\beta,0}$
    followed by the projection of the ancilla onto $\ket{\varphi_F}$ is
    \begin{align}
        \psi(\sigma)
         & =\bra{\varphi_F}A_0^{\sigma_0}A_1^{\sigma_1}\cdots
        A_{N-1}^{\sigma_{N-1}}\ket{\varphi_I}.
        \label{eq:seqanc_state}
    \end{align}
    This is \cite[Section~5.1]{PerezGarcia2007Matrix}.
\end{definition}

\begin{definition}[Probabilistic sequential scheme]
    \label{def:seqanc_probabilistic}
    \lean{MPSPreparation.IsProbabilisticallyGenerated}
    \leanok
    \uses{def:seqanc_state}
    A vector $\psi\in(\C^d)^{\otimes N}$ is \emph{probabilistically generated
        with a $D$-dimensional ancilla} if it has the form
    (\ref{eq:seqanc_state}) for some matrices $A_k^i\in\MN{D}$ and some
    vectors $\ket{\varphi_I},\ket{\varphi_F}\in\C^D$. The vector is the
    unnormalized post-measurement vector, as in the source's ``up to
    normalization''.
\end{definition}

\begin{definition}[Operations induced by unitaries]
    \label{def:seqanc_step}
    \lean{MPSPreparation.stepMatrix, MPSPreparation.jointState}
    \leanok
    Let $d\ge 1$. A unitary $U$ on $\C^D\otimes\C^d$ acting on the site state
    $\ket{0}$ induces the matrices
    $(A^i)_{\alpha\beta}=\bra{\alpha,i}U\ket{\beta,0}$. For unitaries
    $U_0,\dots,U_{N-1}$, with $U_k$ applied at step $k+1$ and acting on site
    $N-1-k$, and an initial ancilla state $\ket{\varphi_I}$, the ancilla
    component of the joint state at the configuration $\sigma$ is
    $A_0^{\sigma_0}\cdots A_{N-1}^{\sigma_{N-1}}\ket{\varphi_I}\in\C^D$, where
    $A_k$ is induced by $U_{N-1-k}$. This is
    \cite[Section~5.1]{PerezGarcia2007Matrix}.
\end{definition}

\begin{definition}[Deterministic sequential scheme]
    \label{def:seqanc_deterministic}
    \lean{MPSPreparation.IsDeterministicallyGenerated}
    \leanok
    \uses{def:seqanc_step}
    Let $d\ge 1$. A vector $\psi$ is \emph{deterministically generated with a
        $D$-dimensional ancilla} if there are unitaries $U_0,\dots,U_{N-1}$ on
    $\C^D\otimes\C^d$ and unit vectors $\ket{\varphi_I},\ket{\varphi_F}\in\C^D$
    such that the ancilla component of the joint state at every
    configuration $\sigma$ equals $\psi(\sigma)\ket{\varphi_F}$; that is, the
    final joint state is $\ket{\varphi_F}\otimes\ket{\psi}$ and the ancilla
    decouples in the last step without measurement.
\end{definition}

\begin{definition}[Open-boundary representation with bounded bond dimension]
    \label{def:seqanc_obc_rep}
    \lean{MPSPreparation.HasOBCRep}
    \leanok
    \uses{def:varying_bond_obc_coeff}
    A vector $\psi$ has an \emph{open-boundary representation with bond
        dimension at most $D$} if $\psi=c_A$ for some varying-bond open chain
    $A$ with common bound $D$, with $c_A$ as in
    (\ref{eq:varying_bond_obc_coeff}).
\end{definition}

\section{Isometric decompositions}

\begin{definition}[Isometry on a bond block]
    \label{def:seqanc_isometry_on}
    \lean{MPSPreparation.IsIsometryOn}
    \leanok
    A family $Q^i\in\MN{D}$, $i=0,\dots,d-1$, is \emph{isometric on its first
        $b$ columns} if for all $\beta,\beta'<b$
    \begin{align}
        \sum_{i=0}^{d-1}\sum_{\alpha=0}^{D-1}
        \overline{Q^i_{\alpha\beta}}\,Q^i_{\alpha\beta'}
         & =\delta_{\beta\beta'}.
        \notag
    \end{align}
    For $b=D$ this is $\sum_i(Q^i)^\dagger Q^i=\Id_D$.
\end{definition}

\begin{definition}[Supports below a level]
    \label{def:seqanc_support}
    \lean{MPSPreparation.IsSupportedBelow, MPSPreparation.IsRowSupportedBelow}
    \leanok
    A vector $v\in\C^D$ is \emph{supported below $b$} if $v_\beta=0$ for
    $\beta\ge b$. A matrix $M\in\MN{D}$ is \emph{row-supported below $a$} if
    $M_{\alpha\beta}=0$ for $\alpha\ge a$.
\end{definition}

\begin{theorem}[Propagation of supports]
    \label{thm:seqanc_support_propagation}
    \lean{MPSPreparation.IsRowSupportedBelow.mul,
        MPSPreparation.IsRowSupportedBelow.mulVec,
        MPSPreparation.isSupportedBelow_eval_mulVec}
    \leanok
    \uses{def:seqanc_support}
    If $M$ is row-supported below $a$, then so is $MN$ for every $N$, and
    $Mv$ is supported below $a$ for every $v$. Consequently, if $Q_p^i$ is
    row-supported below $b_p$ for $p=0,\dots,n-1$ and $v$ is supported
    below $b_n$, then $Q^\sigma v$ is supported below $b_0$.
\end{theorem}
\begin{proof}
    \leanok
    The rows $\alpha\ge a$ of $MN$ and of $Mv$ are sums of multiples of
    vanishing entries. The chain statement follows by induction on $n$,
    applying the vector statement to the first factor.
\end{proof}

\begin{definition}[Boundary vector and boundary matrices]
    \label{def:seqanc_boundary}
    \lean{MPSPreparation.basisVecZero, MPSPreparation.rowMat,
        MPSPreparation.colMat}
    \leanok
    Let $\ket{0}\in\C^D$ be the first standard basis vector. For
    $v,r\in\C^D$ let $\ket{0}v^{\mathsf T}$ be the matrix with first row $v$
    and all other rows zero, and $r\bra{0}$ the matrix with first column $r$
    and all other columns zero.
\end{definition}

\begin{theorem}[Coordinates of the boundary vector and matrices]
    \label{thm:seqanc_boundary_coords}
    \lean{MPSPreparation.isRowSupportedBelow_rowMat,
        MPSPreparation.basisVecZero_dotProduct, MPSPreparation.star_basisVecZero,
        MPSPreparation.mulVec_basisVecZero_apply,
        MPSPreparation.rowMat_mul_mulVec_zero, MPSPreparation.mul_colMat_apply,
        MPSPreparation.mul_colMat_apply_of_ne,
        MPSPreparation.eq_smul_basisVecZero_of_isSupportedBelow_one,
        MPSPreparation.star_smul_basisVecZero_dotProduct_self,
        MPSPreparation.star_basisVecZero_dotProduct_self}
    \leanok
    \uses{def:seqanc_boundary, def:seqanc_support}
    Let $D>0$. The matrix $\ket{0}v^{\mathsf T}$ is row-supported below one;
    $\ket{0}$ is real and of norm one; $\bra{0}w=w_0$ and
    $(M\ket{0})_\alpha=M_{\alpha 0}$; the first coordinate of
    $\ket{0}v^{\mathsf T}Mw$ is $v^{\mathsf T}Mw$; the first column of
    $M\,r\bra{0}$ is $Mr$ and its other columns vanish; a vector supported
    below one equals $w_0\ket{0}$; and $\|x\ket{0}\|^2=|x|^2$.
\end{theorem}
\begin{proof}
    \leanok
    Direct computation of coordinates.
\end{proof}

\begin{theorem}[Peeling off the last site]
    \label{thm:seqanc_eval_last}
    \lean{MPSChainTensor.eval_succ'}
    \leanok
    For square matrices $Q_0,\dots,Q_n$ and a configuration $\sigma$ of $n+1$
    sites, $Q^\sigma=(Q_0^{\sigma_0}\cdots
        Q_{n-1}^{\sigma_{n-1}})\,Q_n^{\sigma_n}$.
\end{theorem}
\begin{proof}
    \leanok
    Associativity of the ordered product.
\end{proof}

\begin{theorem}[One decomposition step]
    \label{thm:seqanc_one_step}
    \lean{MPSPreparation.exists_isIsometryOn_mul}
    \leanok
    \uses{def:seqanc_isometry_on, def:seqanc_support}
    Let $Y^i\in\MN{D}$ be row-supported below $a$ for every $i$. Then there
    are $b\le D$, matrices $Q^i$ and a matrix $R$ with $Y^i=Q^iR$ for every
    $i$, such that each $Q^i$ is row-supported below $a$ and vanishes on the
    columns $\beta\ge b$, the family $Q$ is isometric on its first $b$
    columns, and $R$ is row-supported below $b$. Moreover $b\le c$ whenever
    every $Y^i$ vanishes on the columns $\gamma\ge c$.
\end{theorem}
\begin{proof}
    \leanok
    Stack the matrices $Y^i$ into one matrix with rows indexed by
    $(\alpha,i)$, choose an orthonormal basis $q_0,\dots,q_{b-1}$ of its
    column space, and let $Q^i$ have the columns $q_\beta$ for $\beta<b$ and
    zero columns otherwise. The coefficients of the columns of the stacked
    matrix in this basis form $R$. The column space is spanned by the at most
    $c$ nonzero columns, so $b\le c$. This gives the same factorization as the
    singular value decomposition used in the proof cited in
    \cite[Section~5.1]{PerezGarcia2007Matrix}.
\end{proof}

\begin{theorem}[Successive decompositions]
    \label{thm:seqanc_sweep}
    \lean{MPSPreparation.exists_isometric_chain}
    \leanok
    \uses{def:seqanc_isometry_on, def:seqanc_support}
    Let $a\le D$, let $R_{-1}\in\MN{D}$ be row-supported below $a$, let
    $P_0,\dots,P_{n-1}$ be families of $D\times D$ matrices, and let
    $r\in\C^D$. Then there are $b_0=a$ and $b_1,\dots,b_n\le D$, families
    $Q_0,\dots,Q_{n-1}$ such that $Q_p^i$ is row-supported below $b_p$,
    vanishes on the columns $\beta\ge b_{p+1}$, and $Q_p$ is isometric on its
    first $b_{p+1}$ columns, and a vector $r'$ supported below $b_n$, such
    that for every configuration $\sigma$
    \begin{align}
        R_{-1}\,P^\sigma\,r
         & =Q^\sigma\,r'.
        \notag
    \end{align}
    Moreover $b_{p+1}\le c$ whenever every $P_p^i$ vanishes on the columns
    $\gamma\ge c$.
\end{theorem}
\begin{proof}
    \leanok
    \uses{thm:ldp_sweep_matrix_remainder, thm:seqanc_support_propagation}
    Theorem~\ref{thm:ldp_sweep_matrix_remainder} gives the bonds $b_p$, the
    families $Q_p$ and a matrix $R$ row-supported below $b_n$ with
    $R_{-1}P^\sigma=Q^\sigma R$ for every $\sigma$. Put $r'=Rr$, which is
    supported below $b_n$ by Theorem~\ref{thm:seqanc_support_propagation}; then
    $R_{-1}P^\sigma r=Q^\sigma Rr=Q^\sigma r'$. The induction behind
    Theorem~\ref{thm:ldp_sweep_matrix_remainder} is the construction by
    successive singular value decompositions of Schön, Solano, Verstraete,
    Cirac and Wolf, on which the proof in
    \cite[Section~5.1]{PerezGarcia2007Matrix} is based.
\end{proof}

\begin{theorem}[Isometric chains preserve the norm]
    \label{thm:seqanc_norm}
    \lean{MPSPreparation.sum_normSq_mulVec_of_isIsometryOn,
        MPSPreparation.sum_normSq_eval_mulVec}
    \leanok
    \uses{def:seqanc_isometry_on, def:seqanc_support}
    If $Q$ is isometric on its first $b$ columns and $v$ is supported below
    $b$, then $\sum_i\|Q^iv\|^2=\|v\|^2$. Consequently, if $Q_p^i$ is
    row-supported below $b_p$, $Q_p$ is isometric on its first $b_{p+1}$
    columns, and $v$ is supported below $b_n$, then
    $\sum_\sigma\|Q^\sigma v\|^2=\|v\|^2$.
\end{theorem}
\begin{proof}
    \leanok
    \uses{thm:ldp_sweep_matrix_remainder}
    Take $w=v$ in the inner-product identities of
    Theorem~\ref{thm:ldp_sweep_matrix_remainder}: $\sum_i\langle Q^iv,Q^iv\rangle
        =\langle v,v\rangle$ for one site, and
    (\ref{eq:ldp_sweep_inner_product}) with $w=v$ for the chain.
\end{proof}

\begin{theorem}[Changing sites outside the used block]
    \label{thm:seqanc_block_congr}
    \lean{MPSPreparation.mulVec_eq_of_isSupportedBelow,
        MPSPreparation.eval_mulVec_congr}
    \leanok
    \uses{def:seqanc_support}
    If $M'_{\alpha\beta}=M_{\alpha\beta}$ for all $\beta<b$ and $w$ is
    supported below $b$, then $M'w=Mw$. Consequently, if $Q_p^i$ is
    row-supported below $b_p$, the families $Q'_p$ agree with $Q_p$ on the
    columns $\beta<b_{p+1}$, and $v$ is supported below $b_n$, then
    $Q'^\sigma v=Q^\sigma v$ for every $\sigma$.
\end{theorem}
\begin{proof}
    \leanok
    \uses{thm:seqanc_support_propagation}
    The first claim holds entrywise. The second follows by induction on the
    chain from the right, since each partial product applied to $v$ is
    supported below the next bond level.
\end{proof}

\begin{theorem}[Unitary extension of a block isometry]
    \label{thm:seqanc_unitary_extension}
    \lean{MPSPreparation.exists_unitary_extension}
    \leanok
    \uses{def:seqanc_isometry_on}
    Let $d\ge 1$ and let $Q$ be isometric on its first $b$ columns. Then
    there is a unitary $U$ on $\C^D\otimes\C^d$ with
    $\bra{\alpha,i}U\ket{\beta,0}=Q^i_{\alpha\beta}$ for all $i$, $\alpha$
    and $\beta<b$.
\end{theorem}
\begin{proof}
    \leanok
    The vectors $\sum_{\alpha,i}Q^i_{\alpha\beta}\ket{\alpha,i}$, $\beta<b$,
    are orthonormal. Extend them to an orthonormal basis of
    $\C^D\otimes\C^d$ and let $U$ send the standard basis, with the vectors
    $\ket{\beta,0}$, $\beta<b$, first, to this basis.
\end{proof}

\begin{theorem}[Positive bond bound]
    \label{thm:seqanc_bond_bound_pos}
    \lean{OBCChainTensor.bondBound_pos}
    \leanok
    \uses{def:varying_bond_obc_chain}
    The common bound $D$ of a varying-bond open chain is positive.
\end{theorem}
\begin{proof}
    \leanok
    $1=D_0\le D$.
\end{proof}

\begin{theorem}[Open-chain coefficient as a corner entry]
    \label{thm:seqanc_coeff_corner}
    \lean{OBCChainTensor.coeff_eq_eval_zeroPad}
    \leanok
    \uses{def:varying_bond_obc_coeff, def:varying_bond_obc_zero_pad}
    For a varying-bond open chain with square padding $\widetilde A$, and
    every configuration $\sigma$ of any length $N\ge 0$,
    $c_A(\sigma)=(\widetilde A^\sigma)_{00}$.
\end{theorem}
\begin{proof}
    \leanok
    \uses{thm:varying_bond_obc_coeff_zero, thm:varying_bond_obc_zero_pad_coeff}
    For $N=0$ both sides are one. Otherwise $c_A(\sigma)$ is the trace of
    $\widetilde A^\sigma$, and all diagonal entries except the first vanish
    because the first padded site matrix is row-supported below $D_0=1$.
\end{proof}

\begin{definition}[Cutting a square chain down to its bond blocks]
    \label{def:seqanc_cut_down}
    \lean{OBCChainTensor.ofSupported}
    \leanok
    \uses{def:varying_bond_obc_chain}
    Given $b_0,\dots,b_N\le D$ with $b_0=b_N=1$ and square matrices $Q_k^i$,
    the varying-bond open chain with dimensions $b_k$ has as site matrices
    the upper-left $b_k\times b_{k+1}$ corners of $Q_k^i$.
\end{definition}

\begin{theorem}[Cutting down and padding]
    \label{thm:seqanc_cut_down}
    \lean{OBCChainTensor.zeroPad_ofSupported,
        OBCChainTensor.sum_conjTranspose_mul_ofSupported}
    \leanok
    \uses{def:seqanc_cut_down, def:varying_bond_obc_zero_pad,
        def:seqanc_isometry_on}
    If every $Q_k^i$ is row-supported below $b_k$ and vanishes on the
    columns $\beta\ge b_{k+1}$, then padding the cut-down chain returns
    $Q$. If $Q_k^i$ is row-supported below $b_k$ and $Q_k$ is isometric on
    its first $b_{k+1}$ columns, then the cut-down site matrices $A_k^i$
    satisfy $\sum_i(A_k^i)^\dagger A_k^i=\Id_{b_{k+1}}$.
\end{theorem}
\begin{proof}
    \leanok
    Entrywise comparison; for the second claim the sum over rows $\alpha$
    may be restricted to $\alpha<b_k$.
\end{proof}

\begin{theorem}[Assembling an open chain]
    \label{thm:seqanc_assemble}
    \lean{OBCChainTensor.exists_of_isometric_chain}
    \leanok
    \uses{def:varying_bond_obc_coeff, def:seqanc_isometry_on, def:seqanc_support}
    Let $b_0,\dots,b_{n+1}\le D$ with $b_0=1$, and let $Q_0,\dots,Q_n$ be
    families with $Q_p^i$ row-supported below $b_p$, vanishing on the columns
    $\beta\ge b_{p+1}$, and isometric on the first $b_{p+1}$ columns. For
    every $r'\in\C^D$ there is a varying-bond open chain $A$ with common
    bound $D$ such that $c_A(\sigma)=(Q^\sigma r')_0$ for every $\sigma$,
    $D_k=b_k$ for $k\le n$, and $\sum_i(A_p^i)^\dagger A_p^i=\Id$ for
    $p<n$. If moreover $r'$ is supported below $b_{n+1}$ and $\|r'\|=1$,
    then also $\sum_i(A_n^i)^\dagger A_n^i=1$.
\end{theorem}
\begin{proof}
    \leanok
    \uses{def:seqanc_cut_down, thm:seqanc_cut_down, thm:seqanc_coeff_corner,
        thm:seqanc_eval_last, thm:seqanc_norm, thm:seqanc_boundary_coords}
    Replace the last family by $Q_n^i\,r'\bra{0}$, set the last bond
    dimension to one, and cut down. By Theorems~\ref{thm:seqanc_coeff_corner}
    and~\ref{thm:seqanc_eval_last}, the coefficient at $\sigma$ is the first
    coordinate of $Q^\sigma r'$. If $r'$ is a unit vector supported below $b_{n+1}$, the isometry of
    $Q_n$ on its first $b_{n+1}$ columns gives
    $\sum_i\|Q_n^ir'\|^2=\|r'\|^2=1$.
\end{proof}

\begin{theorem}[Successive decompositions of an open chain]
    \label{thm:seqanc_sweep_obc}
    \lean{OBCChainTensor.exists_isometric_chain_coeff}
    \leanok
    \uses{def:varying_bond_obc_coeff, def:seqanc_isometry_on}
    Let $A$ be a varying-bond open chain of length $n+1$ with dimensions
    $D_k$ and common bound $D$, and let $c\in\C$. There are $b_0=1$ and
    $b_k\le D_k$, families $Q_0,\dots,Q_n$ as in
    Theorem~\ref{thm:seqanc_sweep}, and a vector $r'$ supported below
    $b_{n+1}$, such that $Q^\sigma r'=c\,c_A(\sigma)\ket{0}$ for every
    $\sigma$ and $\|r'\|^2=\sum_\sigma|c\,c_A(\sigma)|^2$.
\end{theorem}
\begin{proof}
    \leanok
    \uses{thm:seqanc_sweep, thm:seqanc_coeff_corner, thm:seqanc_norm,
        thm:seqanc_support_propagation, thm:seqanc_boundary_coords}
    Apply Theorem~\ref{thm:seqanc_sweep} with $a=1$, $R_{-1}=c\ketbra{0}{0}$,
    the square padding of $A$, and $r=\ket{0}$. The padded site matrix
    $\widetilde A_p$ vanishes on the columns $\gamma\ge D_{p+1}$, so
    $b_{p+1}\le D_{p+1}$. The vector $Q^\sigma r'$ is supported below
    $b_0=1$, so it is its first coordinate times $\ket{0}$, and that
    coordinate is $c\,c_A(\sigma)$ by Theorem~\ref{thm:seqanc_coeff_corner}.
    The norm identity is Theorem~\ref{thm:seqanc_norm}.
\end{proof}

\begin{theorem}[Left-canonical open-boundary representation]
    \label{thm:seqanc_left_canonical}
    \lean{OBCChainTensor.exists_isometric_coeff_eq}
    \leanok
    \uses{def:varying_bond_obc_coeff}
    Every varying-bond open chain $A$ of length $n+1$ has a varying-bond
    open chain $A'$ with the same common bound, $c_{A'}=c_A$, and
    $D'_k\le D_k$ for every $k$, such that
    $\sum_i(A'^i_p)^\dagger A'^i_p=\Id$ for every site $p<n$.
\end{theorem}
\begin{proof}
    \leanok
    \uses{thm:seqanc_sweep_obc, thm:seqanc_assemble}
    Apply Theorem~\ref{thm:seqanc_sweep_obc} with $c=1$ and assemble the
    result by Theorem~\ref{thm:seqanc_assemble}. The last bond has
    dimension one on both sides, and the others are bounded by
    Theorem~\ref{thm:seqanc_sweep_obc}.
\end{proof}

\begin{theorem}[Left-canonical representation of a normalized state]
    \label{thm:seqanc_left_canonical_norm}
    \lean{OBCChainTensor.exists_isometric_coeff_eq_of_norm}
    \leanok
    \uses{def:varying_bond_obc_coeff}
    If moreover $\sum_\sigma|c_A(\sigma)|^2=1$, the chain $A'$ of
    Theorem~\ref{thm:seqanc_left_canonical} can be chosen with
    $\sum_i(A'^i_p)^\dagger A'^i_p=\Id$ at every site $p\le n$.
\end{theorem}
\begin{proof}
    \leanok
    \uses{thm:seqanc_sweep_obc, thm:seqanc_assemble}
    In the construction of Theorem~\ref{thm:seqanc_left_canonical}, the
    final vector $r'$ has norm one by Theorem~\ref{thm:seqanc_sweep_obc}, so
    Theorem~\ref{thm:seqanc_assemble} makes the last site isometric as well.
\end{proof}

\section{Sequential generation with an ancilla}

\begin{theorem}[Generated states are open-boundary states]
    \label{thm:seqanc_generated_obc}
    \lean{MPSPreparation.exists_hasOBCRep_of_isProbabilisticallyGenerated}
    \leanok
    \uses{def:seqanc_probabilistic, def:seqanc_obc_rep}
    If a nonzero vector $\psi$ is probabilistically generated with a
    $D$-dimensional ancilla, then $c\,\psi$ has an open-boundary
    representation with bond dimension at most $D$ for some $c\ne 0$.
\end{theorem}
\begin{proof}
    \leanok
    \uses{thm:seqanc_sweep, thm:seqanc_assemble, thm:seqanc_boundary_coords}
    Since $\psi\ne 0$, $D>0$. For $N=0$ take $c=\psi()^{-1}$ and the empty
    chain. For $N\ge 1$ apply Theorem~\ref{thm:seqanc_sweep} with
    $R_{-1}=\ket{0}\bra{\varphi_F}$, the matrices $A_k$, and
    $r=\ket{\varphi_I}$, assemble the result by
    Theorem~\ref{thm:seqanc_assemble}, and take $c=1$.
\end{proof}

\begin{theorem}[Open-boundary states are generated]
    \label{thm:seqanc_obc_generated}
    \lean{MPSPreparation.isProbabilisticallyGenerated_of_hasOBCRep}
    \leanok
    \uses{def:seqanc_probabilistic, def:seqanc_obc_rep}
    If $c\ne 0$ and $c\,\psi$ has an open-boundary representation with bond
    dimension at most $D$, then $\psi$ is probabilistically generated with a
    $D$-dimensional ancilla.
\end{theorem}
\begin{proof}
    \leanok
    \uses{thm:seqanc_coeff_corner, thm:seqanc_bond_bound_pos,
        thm:seqanc_boundary_coords}
    Take the square padding $\widetilde A$ of the representation,
    $\ket{\varphi_F}=\ket{0}$ and $\ket{\varphi_I}=c^{-1}\ket{0}$; then
    $\bra{0}\widetilde A^\sigma c^{-1}\ket{0}=c^{-1}c\,\psi(\sigma)$ by
    Theorem~\ref{thm:seqanc_coeff_corner}.
\end{proof}

\begin{theorem}[Probabilistic sequential generation]
    \label{thm:seqanc_probabilistic}
    \lean{MPSPreparation.isProbabilisticallyGenerated_iff}
    \leanok
    \uses{def:seqanc_probabilistic, def:seqanc_obc_rep}
    A nonzero vector $\psi$ is probabilistically generated with a
    $D$-dimensional ancilla if and only if $c\,\psi$ has an open-boundary
    representation with bond dimension at most $D$ for some $c\ne 0$.
    This is the probabilistic part of \cite[Section~5, Theorem ``Sequential
        generation with ancilla'']{PerezGarcia2007Matrix}.
\end{theorem}
\begin{proof}
    \leanok
    \uses{thm:seqanc_generated_obc, thm:seqanc_obc_generated}
    Theorems~\ref{thm:seqanc_generated_obc}
    and~\ref{thm:seqanc_obc_generated}.
\end{proof}

\begin{theorem}[Unitaries induce isometric operations]
    \label{thm:seqanc_step_isometry}
    \lean{MPSPreparation.isIsometryOn_stepMatrix}
    \leanok
    \uses{def:seqanc_step, def:seqanc_isometry_on}
    The matrices induced by a unitary on $\C^D\otimes\C^d$ satisfy
    $\sum_i(A^i)^\dagger A^i=\Id_D$.
\end{theorem}
\begin{proof}
    \leanok
    This is the $(\ket{\beta,0},\ket{\beta',0})$ entry of $U^\dagger U=\Id$.
\end{proof}

\begin{theorem}[Deterministically generated states are normalized]
    \label{thm:seqanc_deterministic_norm}
    \lean{MPSPreparation.star_dotProduct_self_of_isDeterministicallyGenerated}
    \leanok
    \uses{def:seqanc_deterministic}
    Every deterministically generated vector $\psi$ satisfies
    $\sum_\sigma|\psi(\sigma)|^2=1$.
\end{theorem}
\begin{proof}
    \leanok
    \uses{thm:seqanc_step_isometry, thm:seqanc_norm}
    By Theorems~\ref{thm:seqanc_step_isometry} and~\ref{thm:seqanc_norm},
    $\sum_\sigma\|A^\sigma\varphi_I\|^2=\|\varphi_I\|^2=1$, and
    $\|A^\sigma\varphi_I\|^2=|\psi(\sigma)|^2\|\varphi_F\|^2=|\psi(\sigma)|^2$.
\end{proof}

\begin{theorem}[Deterministic schemes are probabilistic]
    \label{thm:seqanc_deterministic_probabilistic}
    \lean{MPSPreparation.isProbabilisticallyGenerated_of_isDeterministicallyGenerated}
    \leanok
    \uses{def:seqanc_deterministic, def:seqanc_probabilistic}
    Every deterministically generated vector is probabilistically generated
    with an ancilla of the same dimension.
\end{theorem}
\begin{proof}
    \leanok
    Use the induced matrices and project onto $\ket{\varphi_F}$:
    $\bra{\varphi_F}A^\sigma\ket{\varphi_I}
        =\psi(\sigma)\braket{\varphi_F}{\varphi_F}=\psi(\sigma)$.
\end{proof}

\begin{theorem}[Normalized open-boundary states are deterministically generated]
    \label{thm:seqanc_obc_deterministic}
    \lean{MPSPreparation.isDeterministicallyGenerated_of_hasOBCRep}
    \leanok
    \uses{def:seqanc_deterministic, def:seqanc_obc_rep}
    Let $d\ge 1$. If $\sum_\sigma|\psi(\sigma)|^2=1$, $c\ne 0$, and
    $c\,\psi$ has an open-boundary representation with bond dimension at
    most $D$, then $\psi$ is deterministically generated with a
    $D$-dimensional ancilla.
\end{theorem}
\begin{proof}
    \leanok
    \uses{thm:seqanc_sweep_obc, thm:seqanc_unitary_extension,
        thm:seqanc_block_congr, thm:seqanc_boundary_coords}
    For $N=0$ take $\ket{\varphi_I}=\psi()\ket{0}$ and
    $\ket{\varphi_F}=\ket{0}$. Otherwise apply
    Theorem~\ref{thm:seqanc_sweep_obc} with the scalar $c^{-1}$: this gives
    isometric families $Q_p$ and a vector $r'$ with
    $Q^\sigma r'=\psi(\sigma)\ket{0}$ and $\|r'\|^2=\sum_\sigma|\psi(\sigma)|^2=1$.
    Extend each $Q_p$ to a unitary $U_p$ by
    Theorem~\ref{thm:seqanc_unitary_extension}, and take
    $\ket{\varphi_I}=r'$ and $\ket{\varphi_F}=\ket{0}$. The induced matrices
    agree with $Q_p$ on the used columns, so by
    Theorem~\ref{thm:seqanc_block_congr} the ancilla component at $\sigma$
    is $Q^\sigma r'=\psi(\sigma)\ket{0}$. The ancilla decouples because the
    first bond is one-dimensional.
\end{proof}

\begin{theorem}[Deterministic sequential generation]
    \label{thm:seqanc_deterministic}
    \lean{MPSPreparation.isDeterministicallyGenerated_iff}
    \leanok
    \uses{def:seqanc_deterministic, def:seqanc_obc_rep}
    Let $d\ge 1$. A vector $\psi$ is deterministically generated with a
    $D$-dimensional ancilla if and only if $\sum_\sigma|\psi(\sigma)|^2=1$
    and $c\,\psi$ has an open-boundary representation with bond dimension at
    most $D$ for some $c\ne 0$. This is the deterministic part of
    \cite[Section~5, Theorem ``Sequential generation with
        ancilla'']{PerezGarcia2007Matrix}.
\end{theorem}
\begin{proof}
    \leanok
    \uses{thm:seqanc_deterministic_norm,
        thm:seqanc_deterministic_probabilistic, thm:seqanc_generated_obc,
        thm:seqanc_obc_deterministic}
    For the forward direction, $\psi$ is normalized by
    Theorem~\ref{thm:seqanc_deterministic_norm}, hence nonzero, and
    probabilistically generated by
    Theorem~\ref{thm:seqanc_deterministic_probabilistic}; apply
    Theorem~\ref{thm:seqanc_generated_obc}. The converse is
    Theorem~\ref{thm:seqanc_obc_deterministic}.
\end{proof}

\begin{theorem}[Both schemes agree on normalized states]
    \label{thm:seqanc_det_iff_prob}
    \lean{MPSPreparation.isDeterministicallyGenerated_iff_isProbabilisticallyGenerated}
    \leanok
    \uses{def:seqanc_deterministic, def:seqanc_probabilistic}
    Let $d\ge 1$. A vector $\psi$ with $\sum_\sigma|\psi(\sigma)|^2=1$ is
    deterministically generated with a $D$-dimensional ancilla if and only if
    it is probabilistically generated with a $D$-dimensional ancilla.
\end{theorem}
\begin{proof}
    \leanok
    \uses{thm:seqanc_probabilistic, thm:seqanc_deterministic}
    Theorems~\ref{thm:seqanc_probabilistic}
    and~\ref{thm:seqanc_deterministic}.
\end{proof}

\section{The deterministic transition scheme}

\begin{definition}[Deterministic transition scheme]
    \label{def:seqanc_transition}
    \lean{MPSPreparation.IsTransitionInteraction,
        MPSPreparation.transitionStep, MPSPreparation.IsTransitionGenerated}
    \leanok
    \uses{def:seqanc_deterministic}
    Let $d=2$ and let the ancilla be $\C^D\otimes\C^2$, started in a unit
    vector $\ket{\varphi_I}$. In every step an
    arbitrary unitary $W^{[k]}$ acts on the ancilla, followed by the fixed
    interaction $T$ of the ancilla with the current site
    \begin{align}
        \ket{\varphi}\ket{1}\ket{0}
         & \mapsto\ket{\varphi}\ket{0}\ket{1},
         &
        \ket{\varphi}\ket{0}\ket{0}
         & \mapsto\ket{\varphi}\ket{0}\ket{0},
        \label{eq:seqanc_transition}
    \end{align}
    and the ancilla decouples after the last step, that is, the final joint
    state is $\ket{\varphi_F}\otimes\ket{\psi}$ for a unit vector
    $\ket{\varphi_F}\in\C^D\otimes\C^2$. This is the third scheme
    of \cite[Section~5.1]{PerezGarcia2007Matrix}. The step with ancilla
    unitary $W$ leaving the site in $\ket{i}$ is the matrix $S_i(W)$ on
    $\C^D\otimes\C^2$ with entries
    \begin{equation}
        S_i(W)_{x,y}=\delta_{x_2,0}\,W_{(x_1,i),y}.
        \label{eq:seqanc_transition_step}
    \end{equation}
\end{definition}

\begin{theorem}[A transition step on a site in $\ket{0}$]
    \label{thm:seqanc_transition_step}
    \lean{MPSPreparation.transitionStep_eq_of_isTransitionInteraction}
    \leanok
    \uses{def:seqanc_transition}
    Let $T$ be an interaction on $\C^D\otimes\C^2\otimes\C^2$ satisfying
    \eqref{eq:seqanc_transition}, and let $W$ be a matrix on $\C^D\otimes\C^2$.
    Then for all $x,y$ and $i$,
    \begin{equation*}
        \bra{x,i}T(W\otimes\Id)\ket{y,0}=S_i(W)_{x,y}.
    \end{equation*}
\end{theorem}
\begin{proof}
    \leanok
    For $z=(z_1,z_2)$, the relations \eqref{eq:seqanc_transition} applied to
    $\ket{\varphi}=\ket{z_1}$ give $T\ket{z_1,z_2,0}=\ket{z_1,0,z_2}$, so
    \begin{equation*}
        \bra{x,i}T\ket{z,0}=\delta_{x_2,0}\,\delta_{x_1,z_1}\,\delta_{i,z_2}.
    \end{equation*}
    Since $(W\otimes\Id)_{(z,j),(y,0)}=\delta_{j,0}\,W_{z,y}$,
    \begin{equation*}
        \bra{x,i}T(W\otimes\Id)\ket{y,0}
        =\sum_z\bra{x,i}T\ket{z,0}\,W_{z,y}
        =\delta_{x_2,0}\,W_{(x_1,i),y}
        =S_i(W)_{x,y}.
    \end{equation*}
\end{proof}

\begin{theorem}[Transition schemes generate the open-boundary states]
    \label{thm:seqanc_transition}
    \lean{MPSPreparation.isTransitionGenerated_iff}
    \leanok
    \uses{def:seqanc_transition, def:seqanc_obc_rep}
    Let $d=2$. A vector $\psi$ is generated by a deterministic transition
    scheme with ancilla $\C^D\otimes\C^2$ if and only if
    $\sum_\sigma|\psi(\sigma)|^2=1$ and $c\,\psi$ has an open-boundary
    representation with bond dimension at most $D$ for some $c\ne 0$. This is
    the transition part of \cite[Section~5, Theorem ``Sequential generation
        with ancilla'']{PerezGarcia2007Matrix}.
\end{theorem}
\begin{proof}
    \leanok
    \uses{thm:seqanc_deterministic, thm:seqanc_unitary_extension,
        thm:seqanc_step_isometry, thm:seqanc_norm}
    By Theorem~\ref{thm:seqanc_deterministic} it suffices to show that the
    transition scheme and the deterministic scheme with ancilla $\C^D$
    generate the same vectors. A step of the transition scheme factors as
    $E R_i$, where $E\ket{\varphi}=\ket{\varphi}\ket{0}$ and
    $(R_i)_{\alpha,y}=W_{(\alpha,i),y}$, and $R_iE$ is the operation
    $\bra{\alpha,i}W\ket{\beta,0}$ of the unitary $W$ on ancilla $\otimes$
    site. Given a deterministic scheme with unitaries $U^{[k]}$, take
    $W^{[k]}=U^{[k]}$ and the initial and final ancilla states
    $E\ket{\varphi_I}$ and $E\ket{\varphi_F}$; the ancilla components are then
    $E$ applied to those of the deterministic scheme. This is the use of the
    fixed interaction in the work on which
    \cite[Section~5.1]{PerezGarcia2007Matrix} is based.

    Conversely, after the first step the tag qubit is $\ket{0}$, so the
    ancilla component at $\sigma$ is $E$ applied to the product of the
    matrices $R_iE$ acting on $v_i=R_i\ket{\varphi_I}$, where $R_i$ comes from
    the first ancilla unitary. Since $\sum_i\|v_i\|^2=\|\varphi_I\|^2=1$, the
    vectors $v_i$ are the first column of a unitary on $\C^D\otimes\C^2$ by
    Theorem~\ref{thm:seqanc_unitary_extension}, and the scheme becomes a
    deterministic scheme with ancilla $\C^D$ started in $\ket{0}$. The final
    state is $\psi(\sigma)$ times the tag-zero part of $\varphi_F$. The total
    squared norm is $1$ by Theorem~\ref{thm:seqanc_norm}, so some
    $\psi(\sigma)$ is nonzero, the tag-one part of $\varphi_F$ vanishes, and the
    tag-zero part is normalized.
\end{proof}

\section{Sequential generation without an ancilla}

\begin{definition}[Sequential scheme without an ancilla]
    \label{def:seqanc_no_ancilla}
    \lean{MPSPreparation.noAncillaState,
        MPSPreparation.IsProbabilisticallyGeneratedWithoutAncilla,
        MPSPreparation.IsDeterministicallyGeneratedWithoutAncilla}
    \leanok
    Let $d\ge 1$ and $N\ge 2$. Starting from $\ket{0}^{\otimes N}$, an operation $U^{[k]}$
    acts on sites $k$ and $k+1$, for $k=1,\dots,N-1$ in turn. The scheme is
    deterministic if every $U^{[k]}$ is unitary and probabilistic if the
    operations are arbitrary \cite[Section~5.2]{PerezGarcia2007Matrix}. Here
    sites carry the source's labels $1,\dots,N$; source site $j$ is ket
    position $N-j$ in the labelling of this chapter. After
    $U^{[k]}$ the scheme is described by the vector $\chi_k$ on the first
    $k+1$ sites, defined by $\chi_1=U^{[1]}\ket{0,0}$ and
    \begin{equation*}
        \chi_k(\dots,i,\beta)=\sum_\alpha\bra{i,\beta}U^{[k]}\ket{\alpha,0}
        \chi_{k-1}(\dots,\alpha),
    \end{equation*}
    and the generated state is $\chi_{N-1}$.
\end{definition}

\begin{remark}
    The recursion describes the chain state: before $U^{[k]}$ acts, sites
    $k+1,\dots,N$ are still in $\ket{0}$, and site $k$ is not touched
    afterwards, so the state of the whole chain after $U^{[k]}$ is
    $\ket{\chi_k}\otimes\ket{0}^{\otimes(N-k-1)}$. Here the scheme is
    defined by the recursion, and this identification with the action of the
    two-site operations on the whole chain is not proved.
\end{remark}

\begin{theorem}[Sequential generation without an ancilla]
    \label{thm:seqanc_no_ancilla}
    \lean{MPSPreparation.isProbabilisticallyGeneratedWithoutAncilla_iff,
        MPSPreparation.isDeterministicallyGeneratedWithoutAncilla_iff,
        MPSPreparation.noAncillaState_eq_eval}
    \leanok
    \uses{def:seqanc_no_ancilla, def:seqanc_obc_rep}
    Let $d\ge 1$ and $N\ge 2$. The states generated by sequential schemes without an
    ancilla, either deterministically or probabilistically (the latter up to
    normalization), are exactly the states with an open-boundary
    representation of bond dimension at most $d$, normalized in the
    deterministic case. This is \cite[Section~5, Theorem ``Sequential
        generation without ancilla'']{PerezGarcia2007Matrix}, which prints no
    bound on $N$; the scheme presupposes two sites, and the literal
    one-site reading fails \cite{gap:pgvwc07_sequential_no_ancilla_two_sites}.
\end{theorem}
\begin{proof}
    \leanok
    \uses{thm:seqanc_sweep, thm:seqanc_assemble, thm:seqanc_coeff_corner,
        thm:seqanc_norm,
        thm:seqanc_block_congr, thm:seqanc_unitary_extension}
    With $A^{[k]}_{i,\beta\alpha}=\bra{i,\beta}U^{[k]}\ket{\alpha,0}$ the
    recursion for $\chi_k$ gives
    $\psi(\sigma)=\bigl(A^{[N-1]}_{\sigma_1}\cdots
        A^{[1]}_{\sigma_{N-1}}\ket{0}\bigr)_{\sigma_0}$ in ket positions, an
    open-boundary representation with bond dimension at most $d$ by
    Theorems~\ref{thm:seqanc_sweep} and~\ref{thm:seqanc_assemble}. For unitary operations each
    $A^{[k]}$ satisfies $\sum_iA_i^\dagger A_i=\Id$, so the state is
    normalized by Theorem~\ref{thm:seqanc_norm}.

    Conversely, write an open-boundary state as
    $\psi(\sigma)=\bigl(R\,P_1^{\sigma_1}\cdots P_{N-1}^{\sigma_{N-1}}
        \ket{0}\bigr)_{\sigma_0}$ with $R_{j\gamma}$ the entries of the first
    site (Theorem~\ref{thm:seqanc_coeff_corner}), and run the successive
    decompositions of
    Theorem~\ref{thm:seqanc_sweep} starting from $R$: this gives
    $\psi(\sigma)=\bigl(Q_1^{\sigma_1}\cdots Q_{N-1}^{\sigma_{N-1}}
        r'\bigr)_{\sigma_0}$ with every $Q_p$ isometric on its bond block, and
    $\|r'\|=1$ when $\psi$ is normalized. Take $U^{[1]}$ with
    $\bra{i,\beta}U^{[1]}\ket{0,0}=(Q_{N-1}^ir')_\beta$ and, for $k\ge 2$,
    $U^{[k]}$ with $\bra{i,\beta}U^{[k]}\ket{\alpha,0}=(Q_{N-k}^i)_{\beta\alpha}$
    on the used columns; in the deterministic case these are unitary
    extensions by Theorem~\ref{thm:seqanc_unitary_extension}, and
    Theorem~\ref{thm:seqanc_block_congr} identifies the generated state. The
    source argues instead by running the scheme with an ancilla on site $k+1$
    followed by swaps of neighboring sites.
\end{proof}
